\begin{example}
    Las sucesiones $(an^\sigma\log^\tau n)_{n\geq1}$ con $a \neq 0$, 
    $\sigma \in (0,1)$ y $\tau \in \mathbb{R}$ son equidistribuidas módulo $1$.
\end{example}

\begin{resolve}
    Sea $g(x) = ax^\sigma\log^\tau (x)$ en $[1,\infty)$.
    
    La primera condición se cumple,
    \[
        \lim_{x \to \infty} \frac{ax^\sigma\log^\tau (x)}{x} = 
        \lim_{x \to \infty} \frac{a\log^\tau (x)}{x^{1-\sigma}} = 0
    \]
    porque $\log^c (x) \ll x^d$ para cualquier $c \in \mathbb{R}$ y $d > 0$.
    
    Para la segunda condición, calculamos $g'(x)$,
    \[
        g'(x) = 
        a \left(\sigma x^{\sigma-1} \log^\tau (x) + \tau x^{\sigma-1} \log^{\tau-1} (x)\right) =
        a x^{\sigma-1} \log^{\tau-1} (x) \left(\sigma \log (x) + \tau\right)
    \]
    y por tanto
    \[
        \lim_{x \to \infty} x |a x^{\sigma-1} \log^{\tau-1} (x) \, \left(\sigma \log (x) + \tau\right)| =
        \lim_{x \to \infty} |a| x^{\sigma} \log^{\tau-1} (x) \, \left(\sigma \log (x) + \tau\right) = \infty
    \]

    Justamente no son monotonas en todo $[1,\infty)$
    Modificar el teorema de Fejer para que sea monotona a partir de cierto punto.
\end{resolve}